\item \textbf{Diseño de un doblete óptico acromático.}
Dos lentes delgadas hechas de vidrios con índices $n_1$ y $n_2$ se colocan juntas en contacto para formar un doblete acromático. La primera lente tiene una superficie plana y una cóncava de radio de curvatura $R$. La segunda lente tiene dos superficies convexas, ambas de radio de curvatura $R$. Demuestre que este doblete equivale a una sola lente delgada cuya distancia focal $f$ cumple:
$$ \frac{1}{f} = \frac{2n_2 - n_1 - 1}{R} $$